\section{ChatGPT Agent：连接研究与实践}

Kimi K2 展示模型训练管线，本章转到产品系统，核心问题是：训练出来的 Agent 如何进入用户工作流，并承担真实操作风险。第二份材料是 ChatGPT Agent 发布文。OpenAI 将 Operator、Deep Research 和 ChatGPT 对话能力整合成统一的智能体系统。它通过自己的虚拟计算机、文本浏览器、可视化浏览器、终端和 API 访问来完成任务，同时保留用户确认、接管和中断能力。

\begin{figure}[H]
\centering
\includegraphics[width=0.96\textwidth]{chatgpt-agent-unified-system.png}
\caption{ChatGPT Agent 统一系统：Operator + Deep Research + ChatGPT + 工具计算机。自制概念图，依据 OpenAI 官方页面和 00:43:50--01:53:38 对谈内容整理。}
\end{figure}

\begin{knowledgebox}{读图：产品系统比模型多很多层}
ChatGPT Agent 不只是模型，它包含浏览、搜索、代码执行、文件处理、连接器、用户确认和安全策略。这些层决定它能不能处理真实工作流。
\end{knowledgebox}

\begin{knowledgebox}{ChatGPT Agent 工具栈}
\begin{tabularx}{\textwidth}{p{0.22\textwidth}p{0.34\textwidth}X}
\toprule
工具/层 & 作用 & 风险控制 \\
\midrule
视觉浏览器 & 与人类网页交互，点击、筛选、填写 & 用户可接管浏览器。 \\
文本浏览器 & 高效读取大量网页文本 & 需要来源过滤和引用。 \\
Terminal & 运行代码、处理文件和数据 & 高风险操作需限制。 \\
Connectors/API & 连接 Gmail、GitHub、日历等 & 权限和隐私控制。 \\
User confirmation & 执行重要操作前确认 & 降低误操作后果。 \\
\bottomrule
\end{tabularx}
\end{knowledgebox}

\subsection{安全：新能力带来新风险}

OpenAI 页面强调，新型 Agent 能在网页上执行操作，因此带来提示注入、误操作、隐私和高风险任务问题。控制措施包括重要操作确认、接管模式、监控模式、拒绝高风险任务、限制数据访问和安全浏览器接管。

\begin{warningbox}{Agent 安全的核心变化}
聊天模型说错话，后果通常还在文本层；Agent 做错动作，可能影响账户、文件、邮件、购买、代码和真实业务。因此安全从内容审核扩展到权限、工具、确认和审计。
\end{warningbox}

\begin{knowledgebox}{安全风险分层}
\begin{tabularx}{\textwidth}{p{0.20\textwidth}p{0.34\textwidth}X}
\toprule
风险 & 例子 & 防护思路 \\
\midrule
提示注入 & 网页隐藏指令诱导 Agent 泄露数据 & 检测、隔离、重要动作确认。 \\
误操作 & 误发邮件、误购买、误删文件 & 人类确认、接管、权限分级。 \\
隐私泄露 & 连接器或登录网站数据被读取 & 最小权限、会话清理、禁用无关连接器。 \\
高风险任务 & 金融、生物、危险操作 & 拒绝策略、监控模式、专门防护。 \\
\bottomrule
\end{tabularx}
\end{knowledgebox}

\subsection{本章小结}

ChatGPT Agent 的重点是统一产品系统：它把研究能力、网页操作、工具执行和用户控制放进同一工作流，展示了 Agent 从实验到产品的路径。

